
\section*{Equivalence of Derivative Definitions}

\noindent \textbf{Proposition:} A function $f: (a,b) \to \mathbb{R}$ is differentiable at $c \in (a,b)$ if and only if 
\[ \lim_{h \to 0} \frac{f(c+h) - f(c)}{h} = f'(c). \]

\begin{proof}
$(\Rightarrow)$ Assume $f$ is differentiable at $c$, meaning $\lim_{x \to c} \frac{f(x) - f(c)}{x - c} = f'(c)$.
\begin{enumerate}
    \item \textbf{Substitution:} Let $h = x - c$. It follows that $x = c + h$.
    \item \textbf{Change of Limit Variable:} As $x \to c$, the difference $(x - c) \to 0$, therefore $h \to 0$.
    \item \textbf{Apply Substitution to the Limit:}
    \[ \lim_{x \to c} \frac{f(x) - f(c)}{x - c} = \lim_{h \to 0} \frac{f(c + h) - f(c)}{h} \]
    \item \textbf{Conclusion:} Since the first limit exists and equals $f'(c)$, the second limit must also exist and equal $f'(c)$.
\end{enumerate}

\vspace{1em}

$(\Leftarrow)$ Conversely, assume the limit with respect to $h$ exists and equals $f'(c)$.
\begin{enumerate}
    \item \textbf{Reverse Substitution:} Let $x = c + h$, which implies $h = x - c$.
    \item \textbf{Change of Limit Variable:} As $h \to 0$, $x \to c$.
    \item \textbf{Conclusion:} By substituting back, we recover the original definition:
    \[ \lim_{h \to 0} \frac{f(c+h) - f(c)}{h} = \lim_{x \to c} \frac{f(x) - f(c)}{x - c} = f'(c). \]
    Thus, $f$ is differentiable at $c$.
\end{enumerate}
\end{proof}

